\section{Tổng quan Final Submission}

Final Submission là một tài liệu hợp nhất, không phải Phase 4 tạo thêm một baseline mới. Báo cáo phải phản ánh phiên bản cuối cùng và nhất quán của Submission 1, 2 và 3; bổ sung bằng chứng MVP hoạt động, chuỗi màn hình demonstration và khai báo rõ việc sử dụng Generative AI.

\subsection{Mục tiêu và kết quả đạt được}

Tóm tắt vấn đề, phạm vi hệ thống, các actor chính, bảy phân hệ, MVP đã hiện thực và giá trị có thể quan sát được. Chỉ tuyên bố chức năng hoặc chỉ số đã có bằng chứng.

\subsection{Phạm vi hiện thực cuối cùng}

\begin{longtable}{|L{25mm}|L{42mm}|L{38mm}|L{30mm}|}
  \caption{Phạm vi hiện thực cuối cùng theo phân hệ}\\
  \hline
  \rowcolor{gray!15}
  \textbf{Phân hệ} & \textbf{Use case/feature đã hiện thực} & \textbf{Bằng chứng} & \textbf{Giới hạn còn lại}\\\hline
  \endfirsthead
  \hline
  \rowcolor{gray!15}
  \textbf{Phân hệ} & \textbf{Use case/feature đã hiện thực} & \textbf{Bằng chứng} & \textbf{Giới hạn còn lại}\\\hline
  \endhead
  \texttt{SUB-*} & UC/feature thực sự chạy được & Screen, test, log hoặc demo step & Chưa làm hoặc đang mô phỏng\\\hline
\end{longtable}

\subsection{Cách đọc báo cáo}

Báo cáo không lặp lại ba submission theo thứ tự thời gian. Phần đầu khóa baseline cấp hệ thống; phần tiếp theo gom toàn bộ artefact theo bảy phân hệ và 18 use case để người đọc theo dõi liền mạch từ yêu cầu đến UI, hành vi, class/method, test và demonstration. Deployment view, implementation view và các quyết định dùng chung được trình bày một lần ở phần kiến trúc tổng thể.
